\section{Abelian presheaves}
\label{section-abelian-presheaves}

\noindent
In this section we briefly point out some features of the
category of presheaves that allow one to define presheaves
of abelian groups.

\begin{example}
\label{example-singleton-presheaf}
Let $X$ be a topological space. Consider a rule $\mathcal{F}$ that
associates to every open subset of $X$ a singleton set. Since every set
has a unique map into a singleton set, there exist unique restriction
maps $\rho^U_V$. The resulting structure is a presheaf of sets on $X$.
It is a final object in the category of presheaves of sets on $X$, by the
property of singleton sets mentioned above. Hence it is
also unique up to unique isomorphism. We will sometimes write $*$
for this presheaf.
\end{example}

\begin{lemma}
\label{lemma-product-presheaves}
Let $X$ be a topological space. The category of presheaves of sets
on $X$ has products (see
Categories, Definition \ref{categories-definition-product}).
Moreover, the set of
sections of the product $\mathcal{F} \times \mathcal{G}$
over an open $U$ is the product of the sets of sections of
$\mathcal{F}$ and $\mathcal{G}$ over $U$.
\end{lemma}

\begin{proof}
Namely, suppose $\mathcal{F}$ and $\mathcal{G}$ are
presheaves of sets on the topological space $X$.
Consider the rule $U \mapsto \mathcal{F}(U) \times \mathcal{G}(U)$,
denoted $\mathcal{F} \times \mathcal{G}$. If $V \subset U \subset X$
are open then define the restriction mapping
$$
(\mathcal{F} \times \mathcal{G})(U)
\longrightarrow
(\mathcal{F} \times \mathcal{G})(V)
$$
by mapping $(s, t) \mapsto (s|_V, t|_V)$. Then it is immediately
clear that $\mathcal{F} \times \mathcal{G}$ is a presheaf.
Also, there are projection maps
$p : \mathcal{F} \times \mathcal{G} \to \mathcal{F}$
and
$q : \mathcal{F} \times \mathcal{G} \to \mathcal{G}$.
We leave it to the reader to show that
for any third presheaf $\mathcal{H}$ we have
$\Mor(\mathcal{H}, \mathcal{F} \times \mathcal{G})
= \Mor(\mathcal{H}, \mathcal{F}) \times
\Mor(\mathcal{H}, \mathcal{G})$.
\end{proof}

\noindent
Recall that if $(A, + : A \times A \to A, - : A \to A, 0\in A)$
is an abelian group, then the zero and the negation maps are
uniquely determined by the addition law. In other words, it
makes sense to say ``let $(A, +)$ be an abelian group''.

\begin{lemma}
\label{lemma-abelian-presheaves}
Let $X$ be a topological space.
Let $\mathcal{F}$ be a presheaf of sets.
Consider the following types of structure on $\mathcal{F}$:
\begin{enumerate}
\item For every open $U$ the structure of an abelian group
on $\mathcal{F}(U)$ such that all restriction maps are
abelian group homomorphisms.
\item A map of presheaves
$+ : \mathcal{F} \times \mathcal{F} \to \mathcal{F}$,
a map of presheaves $- : \mathcal{F} \to \mathcal{F}$
and a map $0 : * \to \mathcal{F}$
(see Example \ref{example-singleton-presheaf})
satisfying all the axioms of $+, -, 0$ in a usual
abelian group.
\item A map of presheaves
$+ : \mathcal{F} \times \mathcal{F} \to \mathcal{F}$,
a map of presheaves $- : \mathcal{F} \to \mathcal{F}$
and a map $0 : * \to \mathcal{F}$
such that for each open $U \subset X$ the quadruple
$(\mathcal{F}(U), +, -, 0)$ is an abelian group,
\item A map of presheaves $+ : \mathcal{F} \times \mathcal{F}
\to \mathcal{F}$ such that for every open $U \subset X$
the map $+ : \mathcal{F}(U) \times \mathcal{F}(U) \to \mathcal{F}(U)$
defines the structure of an abelian group.
\end{enumerate}
There are natural bijections between the collections of
types of data (1) - (4) above.
\end{lemma}

\begin{proof}
Omitted.
\end{proof}

\noindent
The lemma says that to give an abelian group object $\mathcal{F}$
in the category of presheaves is the same as giving a presheaf of sets
$\mathcal{F}$ such that all the sets $\mathcal{F}(U)$ are endowed
with the structure of an abelian group and such that all the restriction
mappings are group homomorphisms. For most algebra structures
we will take this approach to (pre)sheaves of such objects, i.e.,
we will define a (pre)sheaf of such objects to be a (pre)sheaf
$\mathcal{F}$ of sets all of whose sets of sections $\mathcal{F}(U)$
are endowed with this structure compatibly with the restriction
mappings.

\begin{definition}
\label{definition-abelian-presheaves}
Let $X$ be a topological space.
\begin{enumerate}
\item A {\it presheaf of abelian groups on $X$} or an
{\it abelian presheaf over $X$}
is a presheaf of sets $\mathcal{F}$ such that for each open
$U \subset X$ the set $\mathcal{F}(U)$ is endowed with
the structure of an abelian group, and such that all restriction
maps $\rho^U_V$ are homomorphisms of abelian groups, see
Lemma \ref{lemma-abelian-presheaves} above.
\item A {\it morphism of abelian presheaves over $X$}
$\varphi : \mathcal{F} \to \mathcal{G}$ is a morphism of presheaves
of sets which induces
a homomorphism of abelian groups $\mathcal{F}(U) \to \mathcal{G}(U)$
for every open $U \subset X$.
\item The category of presheaves of abelian groups on $X$ is denoted
$\textit{PAb}(X)$.
\end{enumerate}
\end{definition}

\begin{example}
\label{example-direct-sum-points}
Let $X$ be a topological space. For each $x \in X$ suppose
given an abelian group $M_x$. For $U \subset X$ open
we set
$$
\mathcal{F}(U) = \bigoplus\nolimits_{x \in U} M_x.
$$
We denote a typical element in this abelian group by
$\sum_{i = 1}^n m_{x_i}$, where $x_i \in U$ and $m_{x_i} \in M_{x_i}$.
(Of course we may always choose our representation such that
$x_1, \ldots, x_n$ are pairwise distinct.)
We define for $V \subset U \subset X$ open a restriction
mapping $\mathcal{F}(U) \to \mathcal{F}(V)$ by
mapping an element $s = \sum_{i = 1}^n m_{x_i}$
to the element $s|_V = \sum_{x_i \in V} m_{x_i}$.
We leave it to the reader to verify that this is a
presheaf of abelian groups.
\end{example}



